Could the intelligence transition become an extraordinary opportunity, and when does anticipating it help an agent build something useful? We define \emph{constructive foresight} as a policy that acquires decision-relevant information and prepares a qualified capability before its delivery window closes. \emph{Mission Alpha} is incremental beneficiary value after full costs relative to a credible alternative. The study separates observed numerical capability, prospective local timing forecasts, modeled adoption and unobserved commercial acceptance. A preregistered instrument on four fixed public sparse matrices passes all 432 held-out batches and 3,024 manufactured right-hand sides. Reuse is 7.0--13.6 times faster than repeated setup at 16 right-hand sides, but conventional cached factorization explains the benefit: added certification has median time ratios of 1.01--1.44 relative to that strong baseline. Frozen primary-deadline forecasts achieve Brier \PrimaryBrier{} versus \BaseBrier{} for a pooled constant, on 720 correlated local outcomes. In 4,608 paired simulated worlds, constructive foresight's nominal increment over the calibration-selected rule is $+0.0257$ assumed credits per episode (95\% seed-block interval $[-0.0389,0.1273]$); reversed signals produce a loss. A compiled conventional policy with identical information chooses the same actions throughout and costs less to run. Published Alpha, Jobs and Manager components demonstrate knowledge-only restoration, fresh proof and separate explicit authority, with 14 expected refusals retained. The result is a working readiness method and a precise negative boundary: this study does not establish a unique full-cost advantage for foresight or observed customer value. Readiness is useful only when it pays for itself against the best available response.
